\chapter{Considerações Finais}
\label{chap:consideracoes_finais}

Esta dissertação estudou o prêmio de risco de volatilidade do Bitcoin, medido
pela diferença entre a volatilidade histórica dos 30 dias seguintes e a
volatilidade implícita dada pelo DVOL, com dados diários de março de 2021 a
março de 2026. O prêmio existe e é grande: a média do BVRP é de $-8{,}76$ p.p.
($t = -5{,}15$), e a IV supera a VH dos 30 dias seguintes em 72\% das datas.
As variáveis conhecidas em $t$ contêm alguma informação sobre o prêmio: na
janela expansiva, todos os modelos de previsão têm erro fora da amostra menor
que o da média histórica, o Ridge reduz o da média em 22\%, e o teste de
Clark--West rejeita a 5\% para os cinco. A previsibilidade, porém, é modesta: o
\textit{model confidence set} não separa os modelos da média histórica, e as
previsões não antecipam os episódios de prêmio positivo. A medida também muda a leitura dos regimes: com a proxy
retrospectiva, o prêmio é maior no regime de baixa volatilidade, como em
\citet{almeida2024cryptoVRP}; com o BVRP prospectivo, a diferença se inverte.

O BVRP previsto não prevê os retornos do Bitcoin nesta amostra. No modelo
linear, o coeficiente é negativo nos seis horizontes, na direção de
\citet{bollerslev2009expected}, mas nenhum horizonte é significante depois das
correções para o regressor gerado e para as comparações múltiplas, e o teste
conjunto não rejeita ($p = 0{,}37$). A floresta aleatória, que escolhe sozinha
os cortes nas volatilidades, também não encontra a relação: o coeficiente fica
perto de zero em todos os horizontes, e o teste conjunto não rejeita
($p = 0{,}96$). Com cerca de 33 observações independentes fora da amostra, os
testes têm pouco poder, e a ausência de significância não é evidência de
ausência de relação.

O Apêndice~\ref{app:estrategias} estende a análise a regras simples de negociação com o BVRP previsto, em caráter exploratório: nenhuma delas supera o compra e manutenção em índice de Sharpe no mesmo período, e a menor perda máxima que apresentam decorre da menor exposição ao ativo, não de informação sobre a direção dos retornos.

\section{Limitações}
\label{sec:concl-limitacoes}

\textbf{A proxy de passeio aleatório.} A expectativa física da volatilidade
não é observável. O BVRP mede a realização do prêmio, e a proxy, que aproxima
sua expectativa, supõe que a volatilidade histórica siga um passeio aleatório
sem deriva. A proxy aproxima mal a realização do prêmio --- a correlação entre
as duas na mesma data é de $-0{,}02$ ---, e os resultados que a usam, como a
diferença do prêmio entre os regimes, dependem dessa hipótese. O BVRP
previsto, que estima a expectativa com mais variáveis, também a aproxima só de
forma modesta.

\textbf{Dados diários.} A dissertação calcula a volatilidade histórica com
retornos diários. Sem dados intradiários, não é possível estimar a volatilidade
realizada com a precisão da literatura de alta frequência
\citep{andersen2003realized} nem separar os movimentos contínuos dos saltos
nos episódios de estresse, justamente aqueles em que o BVRP fica positivo.

\textbf{Período fora da amostra curto.} As 987 previsões cobrem de junho de
2023 a março de 2026 e, com alvos de 30 dias sobrepostos, equivalem a cerca de
33 observações independentes. Os testes de previsibilidade dos retornos têm,
por isso, pouco poder, e uma relação fraca entre o prêmio e os retornos
poderia passar sem ser detectada.

\textbf{Mudança estrutural com os ETFs.} A amostra inclui a aprovação dos ETFs
de Bitcoin à vista pela SEC, em janeiro de 2024, que abriu um canal regulado
de exposição ao ativo e pode ter mudado a demanda por proteção com opções. A
análise trata o período como um bloco único e não testa a estabilidade dos
parâmetros antes e depois da aprovação. A bimodalidade da VH de 30 dias,
que define os regimes, também perde nitidez a partir de 2025.

\textbf{Concentração na Deribit.} A Deribit calcula o DVOL só com as suas
próprias opções. Como a negociação de opções de Bitcoin se concentra nela, o índice é
a medida natural da volatilidade implícita do ativo, mas o prêmio medido é o
dessa bolsa, e a dissertação não o compara com o de opções negociadas em
outras plataformas.

\section{Pesquisa futura}
\label{sec:concl-pesquisa-futura}

Duas extensões decorrem dessas limitações. Com dados intradiários, a
volatilidade realizada poderia ser estimada com mais precisão e separada em
componentes contínuo e de saltos, e a expectativa física poderia vir de modelos
de previsão da variância, como em \citet{bekaert2014vix}, no lugar do passeio
aleatório. À medida que o período posterior aos ETFs se alonga, a estabilidade
do prêmio e da sua relação com os retornos antes e depois da aprovação poderá
ser testada com mais observações independentes.
